\subsection{GROUPS}

Recall the following definition (§ 2, no. 3, Definition 6).

DEFINITION 1. A set with an associative law of composition, possessing an identity element and under which every element is invertible, is called a group.

In other words, a group is a monoid (§ 2, no. 1, Definition 1) in which every element is invertible. A law of composition on a set which determines a group structure on it is called a group law. If G and H are two groups, a magma homomorphism of G into H is also called a group homomorphism. Such a homomorphism f maps identity element to identity element; for, let e (resp. e') be the identity element of G (resp. H); writing the group laws of G and H multiplicatively, e.e = e, whence \(f(e).f(e) = f(e)\) and, multiplying by \(f(e)^{-1}, f(e) = e'\) . Hence f is unital. It then follows from no. 3 of § 2 that \(f(x^{-1}) = f(x)^{-1}\) for all \(x \in G\) .

Example. In any monoid E the set of invertible elements with the structure induced by that on E is a group. In particular, the set of bijective mappings of a set F onto itself (or set of permutations of F) is a group under the law \((f, g) \mapsto f \circ g\) , called the symmetric group of the set F and denoted by \(S_{F}\) .

In this paragraph, unless otherwise indicated, the law of composition of a group will always be written multiplicatively and e will denote the identity element of such a group law.

A group G is called finite if the underlying set of G is finite; otherwise it is called infinite; the cardinal of a group is called the order of the group.

If a law of composition on G determines a group structure on G, so does the opposite law. The mapping of a group G onto itself which associates with each \(x \in G\) the inverse of x is an isomorphism of G onto the opposite group ( \(\S 2\) , no. 3, Proposition 4).

Following our general conventions (Set Theory, II, § 3, no. 1), we shall denote by \(\mathbf{A}^{-1}\) the image of a subset A of G under the mapping \(x \mapsto x^{-1}\) . But it is important to note that, in spite of the analogy of notation, \(\mathbf{A}^{-1}\) is definitely not the inverse element of A under the law of composition \((\mathbf{X}, \mathbf{Y}) \mapsto \mathbf{XY}\) between subsets of G (recall that XY is the set of xy with \(x \in \mathbf{X}, y \in \mathbf{Y}\) ): the identity element under this law is \(\{e\}\) and the only invertible elements of \(\mathfrak{P}(\mathbf{G})\) under this law are the sets A consisting of a single element (such an A, moreover, certainly has inverse \(\mathbf{A}^{-1}\) ). The identity \((\mathbf{AB})^{-1} = \mathbf{B}^{-1}\mathbf{A}^{-1}\) holds for \(\mathbf{A} \subset \mathbf{G}, \mathbf{B} \subset \mathbf{G}\) . A is called a symmetric subset of \(\mathbf{G}\) if \(\mathbf{A} = \mathbf{A}^{-1}\) . For all \(\mathbf{A} \subset \mathbf{G}, \mathbf{A} \cup \mathbf{A}^{-1}, \mathbf{A} \cap \mathbf{A}^{-1}\) and \(\mathbf{AA}^{-1}\) are symmetric.

